\FrontChapter{Notation and glossary}\label{glossary}
Notation is defined locally when each family is introduced. The same letter
can serve different roles in separate chapters. In particular, $a$ is an
arithmetic weight parameter in Chapter~\ref{ch:arithmetic}, a contraction in
Chapter~\ref{ch:coupled}, and a spectral shift in Appendix~\ref{app:spectral}.

\begin{description}
\setlength{\itemsep}{8pt}
\item[Contracted feedback $T_{q,r}$.] The normalized odd entire solution of
$T'_{q,r}(z)=1+T_{q,r}(qz)T_{q,r}(rz)$, $T_{q,r}(0)=0$, for $0<q,r<1$.
The diagonal notation is $f_q=T_{q,q}$.
\item[Feedback bank $F_\Theta$.] A finite collection of channels with
$h_j\ge0$, $\sum_jh_j=1$ and $0<q_j,r_j<1$, entering
$F_\Theta'=1+\sum_jh_jF_\Theta(q_jz)F_\Theta(r_jz)$.
Zero-weight channels are removed before taking an active-channel maximum.
\item[Maximum modulus and entire order.] $M_f(R)=\max_{|z|=R}|f(z)|$.
For a nonpolynomial entire function,
\[
\rho(f)=\limsup_{R\to\infty}\frac{\log\log M_f(R)}{\log R}.
\]
At positive finite order $\rho$, the lower and upper types used here are the
liminf and limsup of $\log M_f(R)/R^\rho$. Their equality is not assumed.
\item[Coupled pair $(U,V)$.] The normalized entire solution of
$U'=1+U(qz)V(rz)$, $V'=1+U(az)V(bz)$ and $U(0)=V(0)=0$,
with $q,r,a,b\in(0,1)$.
\item[Perron balance.] $\operatorname{spr}P$ denotes the spectral radius of
a matrix. The characteristic exponent satisfies
$\operatorname{spr}\begin{pmatrix}q^\rho&r^\rho\\a^\rho&b^\rho\end{pmatrix}=1$.
A positive right eigenvector gives the proportions of the two type vectors.
\item[Arithmetic depth $A_{a,\beta,Q}$.] The entire sum
$\sum_{n\ge2,\ (n,Q)=1}\mu(n)^2n^{-2a}(1-n^{-\beta})^{2s}$, with
$a>1/2$, $\beta>0$, $Q\ge1$ integer and the real logarithm in the complex power.
The atom $n=1$ is excluded at every depth.
\item[Möbius function and modulus radical.] $\mu(1)=1$.
For a product of $k$ distinct primes, $\mu(n)=(-1)^k$.
If a prime square divides $n$, then $\mu(n)=0$.
The radical $\operatorname{rad}(Q)=\prod_{p\mid Q}p$ records the excluded
prime set, with $\operatorname{rad}(1)=1$.
\item[Squarefree Dirichlet sum.] For $v>1$,
$Z_Q(v)=\zeta(v)\zeta(2v)^{-1}\prod_{p\mid Q}(1+p^{-v})^{-1}$.
The unrestricted notation is $D(v)=Z_1(v)$. Thus
$A_{a,\beta,Q}(0)=Z_Q(2a)-1$.
\item[Prime-sieve measure.] $\nu_{F,\boldsymbol\lambda}$ is the arithmetic
atomic measure obtained by multiplying the mass at $x_n=(1-n^{-2})^2$ by
$\prod_{p\in F}(1+\lambda_p\mathbf1_{p\mid n})$.
It is positive exactly when every real strength satisfies $\lambda_p\ge-1$.
\item[Hankel order and polynomial degree.] For moments $M_j=\int x^j\dd\nu$,
the order-$m$ matrix is $H_m=[M_{i+j}]_{0\le i,j<m}$.
It tests quadratic forms of real polynomials of degree at most $m-1$.
At depth $k$, the matrix is $[M_{k+i+j}]_{0\le i,j<m}$.
An order-five test therefore covers degree four, not degree five.
\item[Signed quadratic form.] For real polynomials,
$Q_\nu(f)=\int f(x)^2\dd\nu(x)$. On a positive measure it is a squared
norm. On a signed measure it may be negative and is not a norm.
\item[Orbit spectral zeta $\mathcal Z_N$.] Let
$d_N(w)=\sum_{n\ge0}(-1)^nw^n/(Nn)!$ and $D_N(z)=d_N(z^N)$, $N\ge2$.
The series $\mathcal Z_N(s)=\sum_jw_j^{-s}$ initially runs over the zeros
$w_j$ of $d_N$, counted with multiplicity, for $\Re s>1/N$.
Each $w_j$ represents one $N$-fold rotation orbit of zeros of $D_N$.
Powers use $\operatorname{Arg}w\in(-\pi,\pi]$.
\item[Rising factorial.] $(s)_k=s(s+1)\cdots(s+k-1)$ for $k\ge1$, and
$(s)_0=1$. It occurs in the shifted spectral residues.
\item[Coefficient conventions.] An ordinary generating function uses
$\sum_na_nz^n$. An exponential generating function uses
$\sum_na_nz^n/n!$. Their product and shift rules differ.
\item[Polyphase section and finite taps.] Polyphase splits a coefficient
sequence by residue class. Finite impulse response (FIR) taps act by finite
ordinary convolution. A phase matrix must retain the carry delay when a tap
crosses a residue-class boundary. This is not the Hurwitz product of
exponential generating functions.
\item[Resolution certificate.] An explicit bound or rational witness for a
structural distinction at a specified radius, precision or polynomial
degree. A sufficient threshold need not be optimal.
\end{description}

The word \emph{packet} describes several analytic, arithmetic and spectral
representations. Its meaning is fixed by the objects and maps in each
construction, not by an assumed equivalence among the families.
